Neural computers point toward a machine form in which a single latent runtime state acts as the computer itself, driving pixels, text, and actions while subsuming what operating systems and interfaces handle today.
In this paper, the main result is that NCs have begun to exhibit early runtime primitives---most notably I/O alignment and short-horizon control---while stable reuse, symbolic reliability, and runtime governance remain unresolved.
Our CNC capability map remains useful as a longer-horizon view, spanning efficiency, computation \& reasoning, memory \& storage, I/O \& control, tool bridges, condition-driven generalization, programmability, and artifact generation.
The map is staged and dependency-informed, but the more immediate gap is still the gap from prototype behavior to usable runtime behavior.
Progress toward CNCs will therefore depend not only on stronger models, but also on whether reuse, consistency, and governance become sustained and testable.
If these gaps continue to close, neural computers will look less like isolated demonstrations and more like a plausible candidate machine form for next-generation computers.

\section*{Acknowledgements}
The authors sincerely thank Yasheng Sun for his early input of the GUI data collection.  
The authors  thank Deyao Zhu and Firas Laakom for their feedback on the manuscript.  
Mingchen Zhuge, Haozhe Liu, Shuming Liu, Wenyi Wang, Wenxuan Zhang, Junjie Fei, and Jürgen Schmidhuber were supported by funding from the King Abdullah University of Science and Technology (KAUST) Center of Excellence for Generative AI (award number 5940) and the SDAIA-KAUST Center of Excellence in Data Science and Artificial Intelligence.
